\documentclass[french]{article}
\usepackage{verbatim}
\usepackage{amsmath,amsfonts,amssymb}
\usepackage{pgfplots}
\usepackage[utf8]{inputenc}
\usepackage[french]{babel}
\usepackage{pgf,tikz,tkz-tab}
\usepackage{mathrsfs}
\usetikzlibrary{arrows}
\usepackage{pstricks}
\pgfplotsset{compat=newest}
\usepackage{geometry}
\geometry{hmargin=2cm,vmargin=2cm}


\begin{document}
%\twocolumn
\begin{center}
\underline{\textbf{Repérage dans le plan et géométrie vectorielle}}
\end{center}
\section{Repérage dans le plan}
\section*{Exercice 1}
\begin{enumerate}


\item Dans le repère $(O,\overrightarrow{i},\overrightarrow{j})$ lire les coordonnées des vecteurs nommés de $\overrightarrow{a}$ à $\overrightarrow{h}$.
\item Parmis les vecteurs de $\overrightarrow{a}$ à $\overrightarrow{h}$, choissisez 3 vecteurs de votre choix et écrire une égalité de la forme $\overrightarrow{z}= x\overrightarrow{i}+y\overrightarrow{j}$ pour chacun des vecteurs choisis.
\item Représenter les vecteurs $\overrightarrow{u} \left( \begin{array}{c}
-3\\
4
\end{array} \right)$, $\overrightarrow{v} \left( \begin{array}{c}
0\\
5
\end{array} \right)$, $\overrightarrow{w} \left( \begin{array}{c}
-2\\
2
\end{array} \right)$.
\item Calculer les coordonnées des vecteurs suivants: $2\overrightarrow{a}, -3\overrightarrow{f}, \overrightarrow{c} + 2\overrightarrow{h}, -\overrightarrow{d}+4\overrightarrow{b}$

\end{enumerate}
\definecolor{xdxdff}{rgb}{0.49019607843137253,0.49019607843137253,1}
\begin{tikzpicture}[line cap=round,line join=round,>=triangle 45,x=1cm,y=1cm]
\begin{axis}[
x=1cm,y=1cm,
axis lines=middle,
grid style=dashed,
ymajorgrids=true,
xmajorgrids=true,
xmin=-8.34871263733203,
xmax=7.027770942492273,
ymin=-2.4895189365633157,
ymax=5.907482643356319,
xtick={-8,-7,...,7},
ytick={-2,-1,...,5},]
\clip(-8.34871263733203,-2.4895189365633157) rectangle (7.027770942492273,5.907482643356319);
\draw [->,line width=0.5pt] (0,0) -- (1,0);
\draw [->,line width=0.5pt] (0,0) -- (0,1);
\draw [->,line width=0.5pt] (1,1) -- (1,3);
\draw [->,line width=0.5pt] (-5,2) -- (-3,4);
\draw [->,line width=0.5pt] (-3,-1) -- (-1,-1);
\draw [->,line width=0.5pt] (-5,-1) -- (-7,1);
\draw [->,line width=0.5pt] (-7,2) -- (-1,1);
\draw [->,line width=0.5pt] (-3,2) -- (2,4);
\draw [->,line width=0.5pt] (2,-1) -- (3,1);
\draw [->,line width=0.5pt] (4,1) -- (3,-1);
\begin{scriptsize}
\draw [fill=xdxdff] (0,0) circle (1pt);
\draw[color=xdxdff] (-0.21599436894058222,-0.11697557170762482) node {$O$};
\draw[color=black] (0.4807525432702065,-0.20496269391088714) node {$\overrightarrow{i}$};
\draw[color=black] (-0.09586559097320485,0.4956811959259967) node {$\overrightarrow{j}$};
\draw[color=black] (1.1294479442940442,2.0333295539084193) node {$\overrightarrow{h}$};
\draw[color=black] (-4.132192530677084,3.222604455785449) node {$\overrightarrow{a}$};
\draw[color=black] (-1.9578616494675538,-1.3140444288368518) node {$\overrightarrow{b}$};
\draw[color=black] (-6.042240100358384,0.415301420509783784) node {$\overrightarrow{c}$};
\draw[color=black] (-3.9519993637260185,1.6969689755997643) node {$\overrightarrow{d}$};
\draw[color=black] (-0.5283291916557634,3.222604455785449) node {$\overrightarrow{e}$};
\draw[color=black] (2.3187228461710805,0.2073721288042924) node {$\overrightarrow{f}$};
\draw[color=black] (3.59208789262528,-0.2130785940815262) node {$\overrightarrow{g}$};
\end{scriptsize}
\end{axis}
\end{tikzpicture}

\section{Coordonnées de vecteurs}

\section*{Exercice 2}
Dans un repère, on se donne les points $A(5;8),B(3;2)$ et $C(2;9)$.

\begin{enumerate}
\item Calculer les coordonnées des vecteur $\overrightarrow{AB}, \overrightarrow{BC}$ et $\overrightarrow{AC}$
\item Calculer les coordonnées du vecteur $\overrightarrow{AB}+\overrightarrow{BC}$
\item Pouvait-on s'attendre à ce résultat ? Justifiez.
\end{enumerate}

\section*{Exercice 3}
Dans un repère, on se donne un vecteur $\overrightarrow{u} \left( \begin{array}{c}
1\\
3
\end{array} \right)$ et un point $A(2;0)$. Déterminer les coordonnées du point $M$ tel que $\overrightarrow{AM}=\overrightarrow{u}$

\section{Déterminant et colinéarité}
\section*{Exercice 4}
À l'aide du déterminant, déterminer si les vecteurs $\overrightarrow{u}$ et $\overrightarrow{v}$ sont colinéaire entre eux.

\begin{enumerate}
\item  $\overrightarrow{u} \left( \begin{array}{c}
-3\\
4
\end{array} \right)$ et $\overrightarrow{v} \left( \begin{array}{c}
0\\
5
\end{array} \right)$

\item  $\overrightarrow{u} \left( \begin{array}{c}
0\\
2
\end{array} \right)$ et $\overrightarrow{v} \left( \begin{array}{c}
0\\
5
\end{array} \right)$



\item  $\overrightarrow{u} \left( \begin{array}{c}
3\\
2
\end{array} \right)$ et $\overrightarrow{v} \left( \begin{array}{c}
\dfrac{15}{2}\\
5
\end{array} \right)$

\item  $\overrightarrow{u} \left( \begin{array}{c}
x\\
y
\end{array} \right)$ et $\overrightarrow{v} \left( \begin{array}{c}
y\\
x
\end{array} \right)$

\end{enumerate}
\section{Longueurs, milieux et géométrie}
\section*{Exercice 5}
On considère les vecteurs  $\overrightarrow{u} \left( \begin{array}{c}
-3\\
2
\end{array} \right)$ et $\overrightarrow{v} \left( \begin{array}{c}
9\\
6
\end{array} \right)$
Démontrer qu'ils sont colinéaires et calculer leurs normes.

\section*{Exercice 6}
Dans le repère $(O,\overrightarrow{i},\overrightarrow{j})$ de l'exercice 1, calculer la norme des vecteurs nommés de $\overrightarrow{a}$ à $\overrightarrow{h}$.

\section*{Exercice 7}
Dans un repère orthonormé, on se donne les points $A(	-3;2)$, $B(3;1)$, $C(5;2)$, et $D(0;-4)$. Calculer la longeur AB, CD et BC.


\section*{Execice 8}
Dans un repère orthonormé, on se donne les points $A(	-3;4)$, $B(0;2)$, $C(4;4)$. Calculer les coordonnées du point K milieu de [AB], L milieu de [BC] et M milieu de [AC].

\section*{Exercice 9}
Dans un repère orthonormé, on considère les points $A(1;0), B(0;2)$ et $C(4;4)$. 
\begin{enumerate}
\item Faire une figure
\item Déterminer la nature du triangle ABC.
\item Déterminer les coordonnées du point D tel que ABCD soit un parallélogramme.
\item Déterminer les coordonnées du centre du parallélogramme ABCD.
\item Calculer la longueur de ses diagonales.
\end{enumerate}

\section*{Exercice 10}
Dans un repère, on se donne les points $A(1;-1)$, $B(0,5;1,5)$, $C(x;4)$. On veut que les points A,B et C soient alignés, notre objectif est de déterminer $x$.
\begin{enumerate}
\item Que doivent respecter les vecteurs $\overrightarrow{AB}$ et $\overrightarrow{BC}$ pour que ABC soient alignés?
\item Déterminer $x$ pour que ABC soient alignés.
\end{enumerate}

\end{document}